\section{Load bytes before validating an animation}
The indexer accepts a complete byte window. A file read, however, can return only
part of what was requested. It can also fail after an earlier successful read.
The next routine, \texttt{video\_load}, establishes the window by reading a
nonempty file through observed end of input. It does not decide whether the
bytes form a valid animation. The caller invokes \texttt{video\_index} only
after loading succeeds.

Pass a NUL-terminated path in \texttt{\$a0}, a staging-buffer address in
\texttt{\$a1}, an unsigned byte capacity in \texttt{\$a2}, and a sixteen-byte
native result context in \texttt{\$a3}. Capacity must be one through
7\,344\,152. The path and byte buffer need no native alignment; the context
is eight-byte aligned. The caller supplies readable path storage, writable
buffer/context storage, and a valid sixteen-byte-aligned stack. These windows
are disjoint. Retain saved registers, the frame pointer and the entry stack
pointer.

\begin{table}[htbp]
  \centering
  \caption{Loader result context. It describes loading, not format validity.}
  \label{tab:video-load-context}
  \begin{tabular}{r p{1.5in} p{2.8in}}
    \toprule
    Offset & Field             & Meaning                                            \\
    \midrule
    0      & Loaded byte count & Eight-byte native output, changed only on success. \\
    8      & Retained suffix   & Eight bytes left unchanged on either return path.  \\
    \bottomrule
  \end{tabular}
\end{table}

Success returns zero in \texttt{\$v0}, the exact positive file count in
\texttt{\$v1}, and that count in context offset zero. Failure returns one and
zero and retains the entire context. The staging buffer may contain a transferred
prefix after failure. Its contents are not an earlier committed file merely
because a previous count remains in the context. If an application must keep an
old animation usable while opening a new one, give the new load separate staging
storage and commit its file, directory and metadata together after validation.
Do not reuse the old animation's backing bytes and then continue using its old
directory after a failed replacement load.

Use the educational file services already introduced in earlier chapters.
Service 13 opens with flags zero for reading and returns a descriptor or a
negative result. Service 14 reads at most its requested byte count, returns a
positive transferred count, zero at EOF, or a negative error. Its count is a
32-bit service argument, so this construction requests at most 4096 bytes at a
time. Addresses and running byte counts remain full-width. Service 16 closes
the owned descriptor and provides no separate close-status result here. Save
the logical result before closing; the service number in \texttt{\$v0} is
not the loader's eventual return status.

When the buffer becomes exactly full, request one byte into private stack storage.
Zero proves EOF at capacity. A positive byte proves overflow; a negative result
is a read error. Neither is success. This probe must not use the first byte beyond
the destination. Every path after a successful open closes its returned descriptor
once, including development paths in unfinished regions. A failed open owns no
descriptor and must not close one.

Attach the loader to your index working file with
\texttt{python3 packages/034-lossless-video/construction/video-io-construction/attach.py --source index-work.asm
  --output reader-work.asm}. The previous index example is retained under
\texttt{index\_example\_main}; the new caller loads and then indexes. Generate
an actual input with \texttt{python3 packages/034-lossless-video/construction/video-io-construction/fixtures.py
  moving --output animation.mlv}. Keep it in the program's working directory or
prepare a path to it. Preparation changes path/capacity data without replacing
instructions and refuses an existing output. The default input arena has the
full conservative capacity; a generated file must still fit the configured
runtime memory window.

Recipes marked for short reads or injected errors use the grading fixture backend;
writing a regular host file alone does not force a particular short-read pattern.
Use ordinary files to check normal loading, missing/empty inputs, exact capacity,
overflow and malformed video content. A nonempty malformed file can load
successfully and then fail indexing. Keep the two statuses distinct when finding
which layer rejected an input.
